\chapter{Einleitung}\label{chap:einleitung}

\section{Motivation und Problemstellung}

\subsection{Fußball als globales Phänomen}

Der Fußball ist die populärste Sportart der Welt: Mit geschätzten 3,5 bis
4 Milliarden Fans weltweit, über 260 Millionen aktiven Spielern und einem
globalen Marktvolumen von mehreren hundert Milliarden Euro jährlich hat
sich das Spiel zu einem kulturellen und ökonomischen Phänomen entwickelt
\cite{Szymanski2003,Statista2024}. Die fünf großen europäischen Ligen
generieren gemeinsam Jahresumsätze von über 25 Milliarden Euro
\cite{DFL2024}. Diese Kommerzialisierung hat den Druck auf Vereine massiv
erhöht: Jede Entscheidung -- von der Kaderzusammenstellung über die Taktik
bis zur Ergebnisvorhersage -- muss zunehmend ökonomischen Kriterien
standhalten \cite{DobsonGoddard2011}.

\subsection{Die zunehmende Rolle von Daten und Analysen im Sport}

Parallel zur Kommerzialisierung hat im Sport eine datengetriebene Revolution
stattgefunden. Während Entscheidungen im Fußball traditionell auf Intuition
und Erfahrung basierten, halten quantitative Methoden zunehmend Einzug
\cite{ReinMemmert2016}. Dieser Wandel firmiert unter Begriffen wie
\glqq Sports Analytics\grqq{} oder \glqq Evidence-Based Coaching\grqq{}
\cite{BunkerThabtah2019}. Die Anwendungsbereiche sind vielfältig:

\begin{itemize}
    \item \textbf{Scouting \& Transfers:} Identifikation unterbewerteter
          Spieler \cite{PappalardoEtAl2019}.
    \item \textbf{Taktikanalyse:} Auswertung von Pressing-Intensität und
          Defensivorganisation \cite{MemmertLemminkSampaio2017}.
    \item \textbf{Verletzungsprävention:} Vorhersage von Risiken aus
          Belastungsdaten \cite{Csercsik2020}.
    \item \textbf{Spielvorhersage:} Prognose von Ergebnissen für strategische
          Planung oder die Wettindustrie \cite{HubacekSourekZelezny2019}.
\end{itemize}

\subsection{Die inhärente Unvorhersehbarkeit des Fußballs}

Trotz aller Fortschritte bleibt eine fundamentale Eigenschaft bestehen:
die hohe Unvorhersehbarkeit von Fußballergebnissen \cite{DobsonGoddard2011}.
In einem Bundesliga-Spiel fallen durchschnittlich lediglich 2,86 Tore
\cite{Transfermarkt2024}. Fußball ist damit ein typisches
\textbf{Low-Scoring-Event}:

\begin{table}[h]
    \centering
    \caption{Vergleich der durchschnittlichen Punktzahlen pro Spiel}
    \label{tab:scoring-comparison}
    \begin{tabular}{l r}
        \toprule
        \textbf{Sportart} & \textbf{Ø Punkte/Spiel} \\
        \midrule
        Fußball (Bundesliga) & 2,86 Tore \\
        Eishockey (NHL)      & 6,2 Tore \\
        American Football (NFL) & 45 Punkte \\
        Basketball (NBA)     & 220 Punkte \\
        \bottomrule
    \end{tabular}
\end{table}

Diese geringe Trefferquote hat weitreichende Konsequenzen. In Sportarten
mit vielen Scoring-Events gleicht sich der Zufall aus -- das bessere Team
setzt sich durch. Im Fußball kann ein einziges Ereignis -- ein abgefälschter
Schuss, ein individueller Fehler -- das Ergebnis entscheidend verzerren
\cite{Sumpter2016}. Die Forschung zeigt, dass selbst ausgereifte Modelle
nur Accuracy-Werte zwischen 50 und 60\,\% erreichen -- geringfügig besser
als eine einfache Baseline \cite{BunkerThabtah2019,HubacekSourekZelezny2019}.
Dennoch: Jede marginale Verbesserung der Vorhersagegenauigkeit kann für
Vereine, Wettanbieter oder Medien erheblichen Wert haben
\cite{ForrestSimmons2000}.

\subsection{Von traditionellen Statistiken zu erweiterten Metriken}

Traditionelle Vorhersagemodelle stützten sich auf Tabellenpositionen, Tore
und Schüsse \cite{DobsonGoddard2011}. Diese Metriken erfassen jedoch nur
oberflächlich, was im Spiel geschieht: Ein Team kann dominieren, 20 Schüsse
abgeben und dennoch verlieren, weil der Gegner mit zwei Torchancen
effizienter war. Die Frage nach der \glqq wahren\grqq{} Teamstärke bleibt
unbeantwortet. Mit \textbf{Expected Goals (xG)} wurde eine Metrik
entwickelt, die die Qualität von Torchancen quantifiziert und damit eine
präzisere Leistungsbewertung ermöglicht \cite{AnzerBauer2021}. xG hat sich
als einer der stabilsten Prädiktoren für zukünftige Teamleistung
etabliert \cite{WheatcroftVanEetvelde2020}.

\section{Expected Goals und Machine Learning in der Spielvorhersage}

xG weist jeder Schusssituation eine Torwahrscheinlichkeit zu, basierend
auf Faktoren wie Schussdistanz, Körperteil und Vorlagentyp
\cite{Spearman2018,AnzerBauer2021}. Das Konzept, seine Berechnung und
verwandte Metriken (xGA, xPTS) werden in
Kapitel~\ref{chap:grundlagen} ausführlich dargelegt.

Die wissenschaftliche Beschäftigung mit Spielvorhersagen reicht von frühen
Poisson-Modellen \cite{Maher1982,DixonColes1997} bis hin zu modernen
Machine-Learning-Verfahren, die nicht-lineare Zusammenhänge erfassen
können \cite{BunkerThabtah2019}. Insbesondere Gradient-Boosting-Methoden
wie XGBoost und LightGBM haben sich für tabellarische Sportdaten als
überlegen erwiesen \cite{GrollLeyEffenberger2019,HubacekSourekZelezny2019}.
Ein detaillierter Überblick über den Forschungsstand findet sich in
Kapitel~\ref{chap:relatedwork}.

\subsection{Forschungslücke: Expected Goals in der Bundesliga}

Trotz dieser Fortschritte bleibt eine zentrale Frage unbeantwortet:
\textbf{Inwieweit verbessert die Integration von Expected-Goals-Daten die
Vorhersagegenauigkeit von Machine-Learning-Modellen für
Bundesliga-Spielergebnisse?}

Während Studien dies für die Premier League \cite{HubacekSourekZelezny2019},
die La Liga \cite{PappalardoEtAl2019} oder europäische Wettbewerbe allgemein
\cite{VanEetveldeEtAl2024} untersuchten, fehlt eine umfassende Analyse für
die deutsche Bundesliga. Dies ist bemerkenswert, da die Bundesliga eigene
Charakteristika aufweist:

\begin{itemize}
    \item \textbf{Hohe Tordichte:} Die Bundesliga hat traditionell eine
          der höchsten Torquoten unter den europäischen Top-Ligen
          \cite{Transfermarkt2024}.
    \item \textbf{Intensives Pressing:} Deutsche Teams sind bekannt für
          aggressives Gegenpressing und hohe Laufleistungen
          \cite{ReinMemmert2016}.
    \item \textbf{Taktische Vielfalt:} Von defensiv ausgerichteten
          Aufsteigern bis zu Possession-Teams der Topklubs.
    \item \textbf{50+1-Regel:} Die besondere Eigentümerstruktur schafft
          ein anderes ökonomisches Umfeld als in England oder Frankreich
          \cite{Mueller2019}.
\end{itemize}

Diese Charakteristika könnten dazu führen, dass Modelle, die für andere
Ligen entwickelt wurden, nicht direkt auf die Bundesliga übertragbar sind.
Eine systematische Validierung für den deutschen Kontext fehlt -- eine
Lücke, die diese Arbeit schließt.

\section{Zielsetzung und Forschungsfrage}

Diese Bachelorarbeit verfolgt das Ziel, die Vorhersagegenauigkeit von
Machine-Learning-Modellen für Bundesliga-Spielergebnisse durch die
systematische Integration von Expected-Goals-Daten zu verbessern und den
Mehrwert dieser Metriken empirisch zu quantifizieren.

Die zentrale Forschungsfrage lautet:

\begin{quote}
    \textbf{Verbessert die Integration von Expected-Goals-Daten (xG) die
    Vorhersagegenauigkeit eines Machine-Learning-Modells für
    Bundesliga-Spielergebnisse signifikant?}
\end{quote}

Diese Hauptfrage lässt sich in drei Teilfragen zerlegen:

\begin{enumerate}
    \item Welche xG-basierten Features sind die stärksten Prädiktoren
          für Spielergebnisse?
    \item Welche Machine-Learning-Algorithmen erreichen die beste
          Vorhersagegenauigkeit?
    \item Wie robust sind die Modelle gegenüber Datenproblemen wie
          Ausreißern und Klassen-Imbalance?
\end{enumerate}

Zur Beantwortung werden zehn Saisons Bundesliga-Spieldaten
(2016/17--2025/26) mit xG-Statistiken von Understat \cite{Understat2024}
verknüpft. Auf Basis von EWA-Features \cite{Brown1956,HyndmanAthanasopoulos2021}
werden fünf Algorithmen (Logistische Regression, Random Forest, LightGBM,
XGBoost, MLP) auf zwölf Datensatzvarianten trainiert und evaluiert. Ein
besonderer methodischer Beitrag ist die konsequente Vermeidung von
Data Leakage durch einen Lag-1-Ansatz sowie speziell konstruierte
Features zur Verbesserung der Unentschieden-Vorhersage \cite{ChawlaEtAl2002}.

\section{Wissenschaftlicher Beitrag und Relevanz}

Diese Arbeit liefert eine der ersten systematischen Evaluationen des
xG-Mehrwerts für die Bundesliga, schließt damit eine in der Literatur
identifizierte Forschungslücke \cite{VanEetveldeEtAl2024} und zeigt, wie
EWA-basiertes Feature Engineering mit Gradient-Boosting-Algorithmen
zu einer reproduzierbaren Pre-Match-Vorhersagepipeline kombiniert werden
kann. Auf praktischer Ebene sind die Ergebnisse für Vereine, die
Wettindustrie \cite{ForrestSimmons2000} und Medien relevant -- überall dort,
wo präzisere Ergebniswahrscheinlichkeiten einen direkten Entscheidungswert
haben.

\section{Aufbau der Arbeit}

Die Arbeit ist wie folgt aufgebaut:

\textbf{Kapitel 2 (Theoretische Grundlagen)} führt in Machine Learning,
die fünf eingesetzten Algorithmen, die Grenzen der Vorhersagbarkeit
multivariater Ereignisse, Dimensionsreduktion sowie das Expected-Goals-Konzept
und Evaluationsmetriken ein.

\textbf{Kapitel 3 (Related Work)} gibt einen Überblick über den Stand der
Forschung zu ML-basierter Spielvorhersage und xG-basierten Modellen.

\textbf{Kapitel 4 (Methodik)} beschreibt das experimentelle Design: den
paarweisen Vergleich mit/ohne xG, die Evaluationsmetriken und die
Data-Leakage-Prävention.

\textbf{Kapitel 5 (Modellierung und Implementierung)} dokumentiert
Datengrundlage, Feature Engineering und die vollständige ML-Pipeline
inkl.\ Hyperparameter-Tuning und Modellspeicherung.

\textbf{Kapitel 6 (Evaluation und Ergebnisvergleich)} präsentiert und
interpretiert die Ergebnisse aller 60 Modellläufe kritisch.

\textbf{Kapitel 7 (Diskussion)} ordnet die Befunde in den wissenschaftlichen
Kontext ein und diskutiert Limitationen.

\textbf{Kapitel 8 (Fazit und Ausblick)} beantwortet die Forschungsfrage
und gibt Empfehlungen für zukünftige Forschung.